\begin{abcliste}
\abc The impedance of a parallel resonant circuit has its peak at the resonance frequency:
\begin{align}
 f_0 &= \foF \\
   &= \frac{1}{2\pi\sqrt{\Lz\times\C}} \\
   &= \fo \approx \result{\foP{3}{4}}
\end{align}
in agreement with the measured peak.

\abc The capacitor stays the same; only the inductance changes:
\begin{align}
 L_1 &= \LaF \\
   &= \frac{1}{\left(2\pi\times\fa\right)^2\times\C} \\
   &= \La \approx \result{\LaP{-6}{3}}
\end{align}

\abc
\begin{align}
 \Delta L &= \dLF \\
   &= \Lz - \La \\
   &= \dL \approx \result{\dLP{-6}{2}}
\end{align}
The frequency rose by about \SI{1}{\percent}, the inductance dropped by about \SI{2}{\percent}: $L$ goes with $1/f^2$, so the frequencies must be measured precisely. A change of a few percent is enough: the detector compares the frequency with a reference and beeps. Iron does the opposite: it raises $L$, and the frequency drops, so a detector can tell iron from coins.
\end{abcliste}
